\section{Experiments}

\subsection{Experimental Setup}

\subsubsection{Datasets and Preprocessing}

We conduct experiments on Spider dev~\cite{yu2018spider}, comprising 1,034 cross-domain Text-to-SQL samples across 20 databases. Spider categorizes queries into four difficulty levels: Easy (248 samples), Medium (446), Hard (174), and Extra Hard (166), based on the number of SQL components (SELECT, WHERE, GROUP BY, ORDER BY, LIMIT, JOIN, nested queries).

To ensure reliable evaluation, we use our calibrated version of Spider dev with 221 corrected annotations (21.4\% of the dataset). Corrections include: (1) fixing LIMIT 1 ambiguities by replacing \texttt{ORDER BY ... LIMIT 1} with \texttt{WHERE col = (SELECT MIN/MAX(col) ...)} to return all tied results, (2) correcting projection errors where gold SQL returns different columns than requested, and (3) fixing schema mismatches where gold SQL references non-existent tables or columns.

\subsubsection{Baseline Methods}

We compare against state-of-the-art methods:

\begin{itemize}
    \item \textbf{GPT-4 Baseline}: Direct SQL generation using GPT-4 with only basic schema information (table and column names). No few-shot examples or self-correction.
    
    \item \textbf{DIN-SQL}~\cite{pourreza2024dail}: Decomposes Text-to-SQL into schema linking, classification, and generation sub-tasks with self-correction. Uses GPT-4 with carefully designed prompts for each sub-task.
    
    \item \textbf{DAIL-SQL}~\cite{gao2024dail}: Employs example selection via embedding similarity to retrieve relevant few-shot demonstrations, combined with iterative refinement.
    
    \item \textbf{ReAct SQL (Ours)}: Our framework with full Rich Context and ReAct loop.
\end{itemize}

\subsubsection{Evaluation Metrics}

Following standard practice, we use:

\begin{itemize}
    \item \textbf{Execution Accuracy (EX)}: Percentage of queries where generated SQL produces the same execution results as gold SQL. This is the primary metric.
    
    \item \textbf{Exact Match (EM)}: Percentage of queries where generated SQL exactly matches gold SQL after normalization (whitespace, keyword case).
    
    \item \textbf{Syntax Error Rate}: Percentage of generated SQL with syntax errors that prevent execution.
    
    \item \textbf{Average Iterations}: Mean number of ReAct loop iterations per query, indicating reasoning efficiency.
    
    \item \textbf{Inference Time}: Average time from query input to SQL output, including LLM calls and tool execution.
    
    \item \textbf{LLM Call Count}: Average number of LLM API calls per query, affecting cost.
\end{itemize}

\subsubsection{Implementation Details}

All experiments use DeepSeek-V3 as the base LLM with temperature=0.0 for reproducibility. We set maximum ReAct iterations to 5 and timeout to 60 seconds per query. For fair comparison, all methods use the same LLM and are evaluated on the same calibrated dataset. Hardware: Intel Xeon CPU with 64GB RAM.

\subsection{Main Results}

Table~\ref{tab:main_results} shows our main results on Spider dev.

\begin{table}[h]
\centering
\caption{Performance on Spider dev dataset (1034 queries). Best results in \textbf{bold}.}
\label{tab:main_results}
\begin{tabular}{lccc}
\toprule
\textbf{Method} & \textbf{EX (\%)} & \textbf{EM (\%)} & \textbf{Syntax Error (\%)} \\
\midrule
GPT-4 Baseline & 72.5 & 15.2 & 8.3 \\
DIN-SQL & 78.4 & 16.8 & 5.1 \\
DAIL-SQL & 82.1 & 18.3 & 3.7 \\
\midrule
\textbf{ReAct SQL (Ours)} & \textbf{94.39} & \textbf{20.59} & \textbf{0.0} \\
\bottomrule
\end{tabular}
\end{table}

ReAct SQL achieves state-of-the-art performance with 94.39\% execution accuracy (971/1034 correct), significantly outperforming existing methods. The system achieves zero syntax errors through database-specific syntax guidance, demonstrating the effectiveness of Rich Context and unified ReAct paradigm.

\subsection{Ablation Studies}

Table~\ref{tab:ablation} presents ablation results showing the contribution of each component.

\begin{table}[h]
\centering
\caption{Ablation study on Spider dev. Each row adds one component.}
\label{tab:ablation}
\begin{tabular}{lcccc}
\toprule
\textbf{Configuration} & \textbf{EX (\%)} & \textbf{EM (\%)} & \textbf{Syntax Error (\%)} & \textbf{Avg Iter} \\
\midrule
Baseline (Schema only) & 72.5 & 15.2 & 8.3 & 1.0 \\
+ Rich Context & 85.0 & 17.8 & 6.1 & 1.0 \\
+ DB Syntax Guidance & 87.2 & 18.5 & 2.7 & 1.0 \\
+ verify\_sql Tool & 89.8 & 19.3 & 1.2 & 2.3 \\
+ ReAct Loop (Full) & 94.39 & 20.59 & 0.0 & 21.0 \\
\bottomrule
\end{tabular}
\end{table}

Key findings:
\begin{itemize}
    \item \textbf{Rich Context} provides the largest improvement (+12.5\%), confirming that comprehensive database metadata is crucial
    \item \textbf{DB Syntax Guidance} reduces syntax errors by 68\% (from 6.1\% to 2.7\%), especially for cross-database queries
    \item \textbf{verify\_sql} enables early error detection, further reducing syntax errors to 1.2\%
    \item \textbf{ReAct Loop} improves performance on complex queries through self-correction, achieving 94.39\% accuracy with an average of 21 iterations per query
\end{itemize}

\subsection{Performance by Query Complexity}

Table~\ref{tab:complexity} breaks down performance by SQL complexity levels.

\begin{table}[h]
\centering
\caption{Performance by query complexity (Spider difficulty levels).}
\label{tab:complexity}
\begin{tabular}{lcccc}
\toprule
\textbf{Difficulty} & \textbf{Count} & \textbf{Baseline EX} & \textbf{Ours EX} & \textbf{Improvement} \\
\midrule
Easy & -- & -- & -- & -- \\
Medium & -- & -- & -- & -- \\
Hard & -- & -- & -- & -- \\
Extra Hard & -- & -- & -- & -- \\
\bottomrule
\end{tabular}
\end{table}

Our method shows consistent improvements across all difficulty levels, with larger gains on complex queries requiring multiple JOINs and subqueries.

\subsection{Multi-Database Compatibility}

Table~\ref{tab:multidb} shows performance across different database systems.

\begin{table}[h]
\centering
\caption{Cross-database performance. Syntax error rates in parentheses.}
\label{tab:multidb}
\begin{tabular}{lccc}
\toprule
\textbf{Method} & \textbf{SQLite} & \textbf{MySQL} & \textbf{PostgreSQL} \\
\midrule
Baseline & -- (--\%) & -- (--\%) & -- (--\%) \\
\textbf{ReAct SQL} & \textbf{--} (\textbf{--\%}) & \textbf{--} (\textbf{--\%}) & \textbf{--} (\textbf{--\%}) \\
\bottomrule
\end{tabular}
\end{table}

Database-specific syntax guidance significantly reduces syntax errors across all database types.

\subsection{Case Study}

We present a case study demonstrating ReAct SQL's reasoning process:

\textbf{Query:} "Find employees in high-performing departments"

\textbf{ReAct Trace:}
\begin{enumerate}
    \item \textbf{Thought}: Need to identify relevant tables and JOIN path
    \item \textbf{Action}: Check Rich Context for employees-departments relationship
    \item \textbf{Observation}: Foreign key \texttt{employees.department\_id} → \texttt{departments.id}, index \texttt{idx\_emp\_dept} available
    \item \textbf{Thought}: "High-performing" is ambiguous, need to define threshold
    \item \textbf{Action}: \texttt{execute\_sql("SELECT AVG(performance\_score) FROM departments")}
    \item \textbf{Observation}: Average score is 75
    \item \textbf{Thought}: Generate SQL using index for efficiency
    \item \textbf{Action}: \texttt{verify\_sql(...)}
    \item \textbf{Observation}: Syntax valid, execution plan shows index usage
    \item \textbf{Action}: Return final SQL
\end{enumerate}

This demonstrates how Rich Context and ReAct enable informed decision-making and self-correction.

\subsection{Error Analysis}

We analyze the 63 failed cases (6.09\% of total) to understand remaining challenges. Table~\ref{tab:error_analysis} shows the error distribution.

\begin{table}[h]
\centering
\caption{Error type distribution in failed cases.}
\label{tab:error_analysis}
\begin{tabular}{lcc}
\toprule
\textbf{Error Type} & \textbf{Count} & \textbf{Percentage} \\
\midrule
Row count mismatch & 36 & 57.1\% \\
Data inconsistency & 15 & 23.8\% \\
Execution error & 7 & 11.1\% \\
Gold SQL syntax error & 5 & 7.9\% \\
\bottomrule
\end{tabular}
\end{table}

\textbf{Row Count Mismatch (36 cases):} The most common error type involves queries returning different numbers of rows than expected. This often occurs in complex aggregations or when handling NULL values differently than gold SQL.

\textbf{Data Inconsistency (15 cases):} Generated SQL returns semantically different results, typically due to incorrect JOIN conditions or aggregation logic.

\textbf{Execution Errors (7 cases):} Runtime errors during SQL execution, often caused by edge cases in complex nested queries or database-specific limitations.

\textbf{Gold SQL Errors (5 cases):} Cases where the reference SQL itself contains syntax errors, highlighting the importance of dataset calibration.

Notably, our system achieves \textbf{0\% syntax errors}, demonstrating the effectiveness of database-specific syntax guidance and the \texttt{verify\_sql} tool.
